% !TeX program = pdflatex
% papers/conversa_arvore_n1_rede_dual.tex — árvore leve curada (rede dual; sem ollama).
\documentclass[11pt,a4paper]{article}
\usepackage[utf8]{inputenc}\usepackage[T1]{fontenc}\usepackage[brazil]{babel}
\usepackage[margin=2.4cm]{geometry}
\usepackage{xcolor}
\usepackage[hidelinks]{hyperref}
\newcommand{\code}[1]{\texttt{\small #1}}
% ─────────────────────────────────────────────────────────────────────────────
%  papers/gkcapa.tex — a capa do Reino, mínima, para os papers em `article`.
%
%  O estilo.tex completo é do `report` (títulos de capítulo, skak, …). Aqui fica
%  só o que a capa precisa, para o paper continuar a compilar sozinho com
%  pdflatex. O tradutor próprio (tests/tex) carrega o estilo.tex da raiz / o
%  symlink papers/estilo.tex e avalia o mesmo `\gkcapa`.
% ─────────────────────────────────────────────────────────────────────────────
\definecolor{ouro}{HTML}{B8912F}
\definecolor{ouroclaro}{HTML}{F3E9CE}
\definecolor{tinta}{HTML}{1E1B16}
\definecolor{regua}{HTML}{6E6858}
\providecommand{\gknota}{\fontsize{7.62}{11.05}\selectfont}
\providecommand{\gktexto}{\fontsize{10.50}{15.22}\selectfont}
\providecommand{\gksub}{\fontsize{12.33}{17.87}\selectfont}
\providecommand{\gksec}{\fontsize{14.47}{20.98}\selectfont}
\providecommand{\gkcap}{\fontsize{16.99}{24.63}\selectfont}
\providecommand{\gktit}{\fontsize{23.42}{33.95}\selectfont}
\providecommand{\Prj}{\mathbb{P}}
\providecommand{\Trf}{\mathcal{T}}
\providecommand{\gkcapa}[3]{%
\title{\vspace{-0.6cm}
{\color{ouro}\rule{\textwidth}{1.4pt}}\\[6mm]
\textbf{\gktit\color{tinta} Reino Dourado}\\[3mm]{\gksec\itshape\color{regua} Golden Kingdom}\\[8mm]
{\gkcap\scshape\color{ouro} #1}\\[5mm]
{\gksub\color{regua} #2}\\[2mm]
{\gktexto\color{regua}\itshape #3}\\[7mm]
{\color{ouro}\rule{\textwidth}{0.6pt}}\\[9mm]{\gknota\color{regua}\begin{minipage}{\textwidth}\setlength{\parindent}{0pt}%
\textbf{Aviso de Isenção Algorítmica:} O universo, as leis e as entidades contidas nestas páginas ---
incluindo felinos matriciais, aves de rapina algébricas e amuletos de controle de fluxo --- são
construções de pura ficção formal. Esta obra não descreve a realidade física, não serve como
engenharia reversa do cosmos e não constitui doutrina religiosa. Qualquer semelhança com bugs reais do seu
sistema operacional é mera coincidência (ou falta de reza). Prossiga por sua própria conta e
risco.\end{minipage}}}
\author{}\date{}}

\begin{document}
\gkcapa{Árvore rede dual}
       {nível 1 curado --- P $\cdot$ $\mathcal{H}$ $\cdot$ H}
       {seguimento de \code{conversa\_rede\_dual.tex}}
\maketitle
\begin{abstract}
Follow-ups curtos sobre a rede neural dual.
Fonte: \code{papers/corpo\_topologico.tex} \S rede-dual;
medidor \code{tests/rede\_dual.js}.
\end{abstract}

\section{falemos de rede\_dual}
Três operadores: percepção $P$, borda $\mathcal{H}$, memória $H$.
Não são três nomes da mesma coisa.
\code{papers/corpo\_topologico.tex} \S rede-dual.

\section{o que a estaca fixa nos pesos}
$W_P=-I$, $b_P=c$ pela Lei~1 --- geometria, não aprendizagem.
A banda $b_{-}-b_{+}=2\Delta$ é memória de retenção.
\code{papers/corpo\_topologico.tex}.

\section{porque hopfield nao inverte sozinho}
Hopfield atrai para $\xi$ (memória $Y$); inversão é $F_H=D\circ F_P^{-1}\circ D^{-1}$.
Atracção $\neq$ reversibilidade.
\code{papers/conversa\_rede\_dual.tex}.

\section{o que fecha lambda mais e menos}
A conjugação: $\lambda_P+\lambda_H=0$ quando $DF_H\,DF_P=I$.
Medido em 1-D e 2-D: \code{tests/rede\_dual.js} §R5--§R7.

\section{falemos de conjugacao}
$D(x)=2c-x$ (estaca Dual Sort) conjuga frente e volta.
UI: \code{app/src/rede\_dual.js}.

\section{o que e X igual a x h}
Estado híbrido $X=(x,h)$; $\mathcal{F}\colon(x,h)\mapsto(x',h')$.
Na banda retém; fora escolhe $F_P$ ou $F_H$.
\code{papers/corpo\_topologico.tex}.

\section{onde a assistente usa a rede}
Ciclo em \code{app/src/assistente.js}: retenção reusa $Y$; volta Hopfield;
frente projecta; Metrónomo atesta $\lambda\Sigma$.

\section{mostra a arvore rede dual}
Este paper + \code{papers/conversa\_rede\_dual.tex} + \code{papers/corpo\_topologico.tex} \S rede-dual.

\end{document}
